An hypotetical shuttle is launched to investigate the atmosphere (100\%
CO$_2$) of two extrasolar planets P$_1$ and P$_2$. The atmosphere is in
static thermodynamic equilibrium. When the shuttle is near each planet, a
radio probe is launched toward respective planet, in vertical direction (in
the direction of the planet's radius). When the radio probe reaches constant
velocity, it starts sending values of the pressure of the atmosphere. In
Fig.~3.1 is plotted the atmospheric pressure values (in arbitrary units) as
function of the time of descent for the planet P$_1$. When the probe touches
the surface of planet P$_1$ it sends the value of the temperature
$T_0 = 700$\,K and the value of the gravitational acceleration
$g_0 = 10\,\mathrm{m\,s}^{-2}$.

% FIGURE NEEDED: 2014_DA_D2-f1.pdf is mis-cropped — the axis tick labels of
% Fig. 3.1 are cut off (y-axis: 20, 40, 60 p-units; x-axis: 1000, 2000,
% 3000 s). Re-crop from 2014_da.pdf page 4.
\ioaafig{2014_DA_D2-f1.pdf}
\begin{center}
Fig.~3.1.
\end{center}

The gravitational acceleration on each planet is assumed to be constant
during uniform descent of the radio probes.

\begin{description}
\item[a)] Find the altitude $h_0$ from where the radio probe R$_1$ starts the
  uniform descent and thus starts the transmitting information.
\item[b)] Find the temperature of planet P$_1$ at the altitude
  $h = 39.6$\,km. You know: The universal gas constant
  $R = 8.3$\,J/(mol\,K); the molar mass of CO$_2$, $\mu = 44$\,g/mol.
\item[c)] In Fig.~3.2.\ was plotted the atmospheric pressure values (in
  arbitrary units) as a function of time of descent for the planet P$_2$
  atmosphere. When the probe touched the surface of the planet P$_2$, it
  sends the value of the temperature $T_0 = 750$\,K and respectively the
  value of gravitational acceleration $g = 8\,\mathrm{m\,s}^{-2}$.

  Draw the following dependency graphs for $p = f(h)$ and $T = f(h)$ in the
  CO$_2$ atmosphere of the planet P$_2$.
\end{description}

% FIGURE NEEDED: 2014_DA_D2-f2.pdf is mis-cropped — the axis tick labels of
% Fig. 3.2 are cut off (y-axis: 20, 40, 60 p-units; x-axis: 500, 1000, 1500,
% 2000 s). Re-crop from 2014_da.pdf page 5.
\ioaafig{2014_DA_D2-f2.pdf}
\begin{center}
Fig.~3.2.
\end{center}
